\documentclass[
  size=12pt, %Possiveis valores 8pt, 9pt, 10pt, 11pt, 12pt, 14pt, 17pt e 20pt
% Aqui você escolhe o estilo, veja páginas 21-25 para
% ver os possíveis estilos.
%  style=upen,    % estilo do powerdot
%  style=tycja,   % estilo do powerdot
%  style=ikeda,   % estilo do powerdot
%  style=simple,
%  style=paintings,
%  style=klope,
%  style=sailor,
%   style=aggie,
%  style=husky,
  style=jefka,
  paper=screen,
%  paper=a4paper,
%% Try me!
  orient=landscape,
%  orient=portrait,
%  mode=print,
  mode=present,
%  mode=handout,
  display=slides,
%  display=slidesnotes,
%  blackslide,
  nopagebreaks,
  hlentries=true, % Realça a tabela de conteúdos da seção atual
  %pauseslide,
  fleqn,
  clock
]{powerdot}

\usepackage[portuges,brazil]{babel}
\usepackage[utf8]{inputenc}
\usepackage{amsmath}
% Aqui vem as fonts especiais
\usepackage{pifont}
\renewcommand\labelitemi{\ding{43}}
\renewcommand\labelitemii{\ding{40}}



%\usepackage{bbold}
%\usepackage{bbm}
\usepackage{textcomp}  % Símbolos matemáticos p/ modo texto
\usepackage{latexsym}  % Símbolos
\usepackage{braket}
\usepackage{fancyhdr}
\usepackage{fancybox}

\usepackage{tcolorbox}
% 
% \usepackage{tikz}
% \usepackage{environ}
% \usetikzlibrary{shapes,snakes}
% 
% 
% \tikzstyle{mybox} = [draw=blue, fill=green!20, very thick,
%     rectangle, rounded corners, inner sep=10pt, inner ysep=20pt]
% 
% \tikzstyle{mynote} = [draw=blue, fill=yellow!50, very thick,
%     rectangle, rounded corners, inner sep=10pt, inner ysep=10pt]
% 
% \tikzstyle{fancytitle} =[fill=yellow!50, text=black, ellipse]
% 
% \NewEnviron{Mybox}{%
% \begin{center}
% \begin{tikzpicture}
% \node [mybox] (box) {%
%     \begin{minipage}[t!]{0.8\textwidth}
%        \BODY
%     \end{minipage}
%     };
% \node[fancytitle] at (box.north) {\textbf{Nota}};
% \end{tikzpicture}
% \end{center}
% }
% 
% \NewEnviron{nota}{%
% \begin{center}
% \begin{tikzpicture}
% \node [mynote] (box) {%
%     \begin{minipage}[t!]{0.8\textwidth}
%        \BODY
%     \end{minipage}
%     };
% \end{tikzpicture}
% \end{center}
% }


%\usepackage[display]{texpower}
%\usepackage{pst-node}
\usepackage{graphicx}
%%\usepackage{rotating}
%\usepackage[dvips]{hyperref}
%\usepackage{fixseminar}
%%\usepackage{color}
%\usepackage{colortbl}
%\usepackage{background}
%\usepackage{pagetrans}
%\usepackage{pp4slide}
%%\usepackage{fancybox}
\usepackage{multicol}
%\usepackage{eepic}
%\usepackage{wrapfig}
%\usepackage{latexcad}
%\usepackage{mpmulti}

\usepackage{listings}
\lstnewenvironment{code}{%
\lstset{frame=single,escapeinside=~~,
backgroundcolor=\color{black!50},
basicstyle=\footnotesize\ttfamily}
}{}

\pdsetup{
  logohook=tl,
%  logopos={.088\slidewidth,.99\slideheight},
  logopos={.0\slidewidth,1.0\slideheight},
  logocmd={\includegraphics[height=0.10\slideheight]{Figs/Simbolo_IF_5.eps}},
%  lf=\textcolor{red}{XXIII Semana da Física},
  lf=\textcolor{red}{XXIII Semana da Física},
%  lf=\textcolor{red}{VI ENCONTRO DE MATEMÁTICA}
  rf=\textcolor{green}{Salviano A. Leão},
  clockformat=h:MM tt,
  clockrefresh=1000,
%  trans=Blinds,
%  trans=Split,
   trans=Box,
%  trans=Wipe,
%   trans=Dissolve,
%  trans=Glitter,
%  trans=Replace,
   theslide=\arabic{slide},
%   palette=River,
%   palette=gold,
%   palette=gray,
%   palette=brown,
%   palette=orange,
%   palette=azul,
%   palette=green,
%   palette=verde,
%   palette=seagreen,  % verde, salviano,
%   palette=GoldenGate,
%   palette=Lamentation,
%   palette=PearlEarring, % GoldenGate, Lamentation, Europa, Moitessier, MayThird, HolyWood, Skater, PearlEarring, Charon,
list={itemsep=6pt}
}


\def\ff#1{\hbox{\rm #1}\,}
\def\fff#1{\hbox{\rm #1}\,}
\def\abs{\ff{abs}}
\def\acos{\ff{acos}}
\def\acosh{\ff{acosh}}
\def\asin{\ff{asin}}
\def\asinh{\ff{asinh}}
\def\asin{\ff{atan}}
\def\asinh{\ff{atanh}}
\def\sin{\ff{sin}}
\def\asen{\ff{asen}}
\def\asenh{\ff{asenh}}
\def\asen{\ff{atg}}
\def\asenh{\ff{atgh}}
\def\sen{\ff{sen}}
\def\cos{\ff{cos}}
\def\tan{\ff{tan}}
\def\sinh{\ff{sinh}}
\def\cosh{\ff{cosh}}
\def\tanh{\ff{tanh}}
\def\log{\fff{log}}
\def\tg{\ff{tg}}
\def\senh{\ff{senh}}
\def\tgh{\ff{tgh}}

\title{Visualização e Tratamento de Dados Com Software Livre}

%\title{Usando Softwares Livres Para Modelar e Tratar Problemas de Física e Ciências Afins}

\author{Salviano A. Leão}
\date{\today}

\begin{document}

\maketitle

%%==============================================================
\begin{slide}[toc=]{Sumário}

\begin{itemize}
   \item Primeiro item \pause
   \item Segundo item  \pause
   \begin{itemize}
      \item Primeiro item do Segundo \pause
      \item Segundo item  \pause
   \end{itemize}

\begin{equation}
   i \hbar \frac{\partial \psi}{\partial t} = H \psi \pause
\end{equation}

   \item Terceiro item \pause
\end{itemize}

\begin{nota}
 Aqui vem um quadro
\end{nota}


\end{slide}


%\overlays{6}{%
% \begin{slide}{Sum�rio}
%   \begin{enumerate}
%     \item Fun��es B�sicas 
%     \item Gr�fico de Fun��es
%     \begin{itemstep}
%        \item Gr�ficos unidimensionais 
%        \item Gr�ficos bidimensionais 
%        \item Gr�ficos de Fun��es Param�tricas
%        \item Gr�ficos em coordenadas polares 
%     \end{itemstep}
%     \item Gr�ficos de Arquivos 
%     \begin{itemstep}
%        \item Gr�ficos unidimensionais 
%        \item Gr�ficos bidimensionais 
%     \end{itemstep}
%     \item Tipos de Terminais 
%     \item Scripts.
%   \end{enumerate}
% \end{slide}
%}

\begin{slide}{Fun��es B�sicas I}
\begin{description}
\item[$\abs(x)$] Para $x$ qualquer, retorna o valor absoluto de $x$ com o mesmo tipo de $x$.
\item[$\abs(x)$] Para $x$ qualquer, retorna $\sqrt{\Re(x)^2+\Im(x)^2}$ para $x$ complexo.
\item[$\acos(x)$] Para $x$ qualquer, retorna um valor entre $-\pi \leq acos(x) \leq \pi$ se $-1\leq x \leq 1$.
\item[$\acosh(x)$] Para $x$ qualquer, retorna um valor entre $-\pi \leq acos(x) \leq \pi$ se $ x \geq 1$.
\item[$arg(x)$] Para $x$ qualquer, retorna a fase de $x$.
\item[$asin(x)$] Para $x$ qualquer, retorna o $\sen^{-1} x$ (inverso do seno).
\item[$asinh(x)$] Para $x$ qualquer, retorna o $\senh^{-1} x$ (inverso do seno hiperb�lico).
\end{description}
\end{slide}

\begin{slide}{Fun��es B�sicas II}
 
\begin{description}
\item[$atan(x)$] Para $x$ qualquer, retorna o $\tg^{-1} x$ (inverso da tangente).
\item[$atan2(y,x)$] Para $x$ e $y$ inteiros ou reais, retorna  $\tg^{-1}(y/x)$ (inverso da tangente).
\item[$atanh(x)$] Para $x$ qualquer, retorna o $\tgh^{-1} x$ (inverso da tangente hiperb�lica).
\item[$besj0(x)$] Para $x$ inteiro ou real, retorna a fun��o de Bessel $j_{0}(x)$.
\item[$besj1(x)$] Para $x$ inteiro ou real, retorna a fun��o de Bessel $j_{1}(x)$.
\item[$besy0(x)$] Para $x$ inteiro ou real, retorna a fun��o de Bessel $y_{0}(x)$.
\item[$besy1(x)$] Para $x$ inteiro ou real, retorna a fun��o de Bessel $y_{1}(x)$.
\end{description}
\end{slide}

\begin{slide}{Fun��es B�sicas III}
 
\begin{description}
\item[$ceil(x)$] Para $\lceil x \rceil$, retorna o menor inteiro n�o menor que $x$ (parte real).
\item[$cos(x)$]  Para $x$ qualquer, retorna o cosseno de $x$, $\cos x$.
\item[$cosh(x)$] Para $x$ qualquer, retorna o cosseno hiperb�lico de $x$, $\cosh x$, .
\item[$erf(x)$] Para $x$ qualquer, retorna a fun��o erro $\hbox{Erf}(\hbox{real}(x))$, do real($x$).
\item[$erfc(x)$] Para $x$ qualquer, retorna $\hbox{Erfc}(\hbox{real}(x)) = 1.0 - \hbox{Erf}(\hbox{real}(x))$.
\item[$exp(x)$] Para $x$ qualquer, retorna $e^{x}$.
\item[$floor(x)$] Para $x$ qualquer, retorna $\lfloor x \rfloor$, o maior inteiro n�o maior do que $x$ (parte real).
\end{description}
\end{slide}




%%==============================================================
\begin{slide}{Sumário}
\tableofcontents
\end{slide}

\section{Parte 1}

%==============================================================
\begin{slide}{Titulo 1}
dafjksda
d

dsafdas
\end{slide}


\section{Parte 2}

%==============================================================
\begin{slide}{Titulo 1}
dafjksda
d

dsafdas
\end{slide}

\section{Parte 3}

%==============================================================
\begin{slide}{Titulo 1}
dafjksda
d

dsafdas
\end{slide}

\end{document}
